% Part: first-order-logic
% Chapter: beyond
% Section: modal-logics

\documentclass[../../../include/open-logic-section]{subfiles}

\begin{document}

\olfileid[hi]{fol}{byd}{mod}

\olsection{विधात्मक तर्कशास्त्र}

किसी सशर्त वाक्य का निम्न उदाहरण देखिए:
\begin{quote}
  यदि जेरेमी उस कमरे में अकेला है, तो वह नशे में है, निर्वस्त्र है और
  कुर्सियों पर नाच रहा है।
\end{quote}
यह ऐसा सशर्त कथन है जो शाब्दिक रूप से सत्य होते हुए भी भ्रामक हो सकता
है, क्योंकि यह केवल इतना कहने के बजाय---कि या तो पूर्वांग असत्य है या
उत्तरांग सत्य---पूर्वांग और निष्कर्ष के बीच किसी अधिक प्रबल संबंध का संकेत
देता है। अर्थात् शब्द-विन्यास बताता है कि दावा केवल इस विशेष जगत में सत्य
नहीं है (जहाँ वह तुच्छ रूप से सत्य हो सकता है, क्योंकि जेरेमी कमरे में
अकेला नहीं है), बल्कि यदि पूर्वांग सत्य \emph{होता}, तो निष्कर्ष भी सत्य
\emph{होता}। दूसरे शब्दों में, कथन का अर्थ लिया जा सकता है कि दावा केवल इस
जगत में नहीं बल्कि हर ``संभावित'' जगत में सत्य है; अथवा कि वह इस विशेष जगत
में मात्र सत्य होने के विपरीत \emph{अनिवार्यतः} सत्य है।

इसी प्रकार की अनिवार्यता को अर्थ देने के लिए विधात्मक तर्कशास्त्र बनाया गया।
सामान्य प्रतिज्ञप्ति-तर्कशास्त्र में एक पेटिका संकारक जोड़कर विधात्मक
प्रतिज्ञप्ति-तर्कशास्त्र मिलता है; अर्थात् यदि $!A$ !!a{formula} है, तो
$\Box !A$ भी है। सहज अर्थ में $\Box !A$ कहता है कि $!A$ \emph{अनिवार्यतः}
सत्य है, अर्थात् हर संभावित जगत में सत्य है। सामान्यतः $\Diamond !A$ को
$\lnot \Box \lnot !A$ का संक्षेप माना जाता है और इसे इस कथन के रूप में पढ़
सकते हैं कि $!A$ \emph{संभवतः} सत्य है। निस्संदेह, विधात्मकता को
विधेय-तर्कशास्त्र में भी जोड़ा जा सकता है।

क्रिप्के !!{structure}s विधात्मक तर्कशास्त्र को अर्थविज्ञान दे सकती हैं;
वस्तुतः क्रिप्के ने पहले इसी तर्कशास्त्र को ध्यान
में रखकर यह अर्थविज्ञान बनाया था। आंशिक क्रमों तक सीमित रहने के बजाय हमारे
पास अधिक सामान्यतः ``संभावित जगतों'' का समुच्चय $P$ और जगतों के बीच एक
द्विपदी ``अभिगम्यता'' संबंध $\Atom{R}{x,y}$ होता है। सहज अर्थ में
$\Atom{R}{p,q}$ कहता है कि जगत~$q$, जगत~$p$ के साथ संगत है; यानी यदि हम
जगत~$p$ ``में'' हैं, तो हमें इस संभावना पर विचार करना होगा कि जगत~$q$ जैसा
हो सकता था।

विधात्मक तर्कशास्त्र को कभी-कभी ``अंतर्वस्त्वात्मक'' तर्कशास्त्र कहा जाता
है, जो ``विस्तारात्मक'' तर्कशास्त्र के विपरीत है। चिरसम्मत तर्कशास्त्र जैसे
विस्तारात्मक तर्कशास्त्र का अभीष्ट अर्थविज्ञान केवल एक, ``वास्तविक'' जगत का
निर्देश करता है; जबकि ``अंतर्वस्त्वात्मक'' तर्कशास्त्र का अर्थविज्ञान अधिक
विस्तृत सत्तामीमांसा पर निर्भर करता है। अनिवार्यता की !!{structure} बनाने
के अतिरिक्त, अनुप्रयोग के अनुसार $\Box$ और $\Diamond$ की पुनर्व्याख्या
करके विधात्मकता से अन्य भाषिक रचनाओं की !!{structure} भी बनाई जा सकती है।
उदाहरणतः:
\begin{enumerate}
\item प्रमाण्यता-तर्कशास्त्र में $\Box !A$ को ``$!A$ प्रमाण्य है'' और
  $\Diamond !A$ को ``$!A$ संगत है'' पढ़ा जाता है।
\item ज्ञानमीमांसात्मक तर्कशास्त्र में $\Box !A$ को ``मैं $!A$ जानता हूँ''
  अथवा ``मैं $!A$ मानता हूँ'' पढ़ सकते हैं।
\item कालिक तर्कशास्त्र में $\Box !A$ को ``$!A$ सदा सत्य है'' और
  $\Diamond !A$ को ``$!A$ कभी-कभी सत्य है'' पढ़ सकते हैं।
\end{enumerate}

हम तर्कशास्त्र में विधात्मकता से संबंधित नियम और अभिगृहीत जोड़ना चाहेंगे।
उदाहरणतः तंत्र $\Log{S4}$ में प्रतिज्ञप्ति-तर्कशास्त्र के सामान्य अभिगृहीतों
और नियमों के साथ निम्न अभिगृहीत होते हैं:
\begin{align*}
& \Box (!A \lif !B) \lif (\Box !A \lif \Box !B)\\
& \Box !A \lif !A \\
& \Box !A \lif \Box \Box !A
\intertext{और यह नियम भी: ``$!A$ से $\Box !A$ निष्कर्षित करो।''
  $\Log{S5}$ निम्न अभिगृहीत जोड़ता है:}
& \Diamond !A \lif \Box \Diamond !A
\end{align*}
इन अभिगृहीतों के रूपांतर अलग-अलग अनुप्रयोगों के अनुकूल हो सकते हैं; जैसे,
सामान्यतः S5 को तार्किक अनिवार्यता की धारणा का अभिलक्षण माना जाता है। एक
अच्छी बात यह है कि यदि क्रिप्के !!{structure}s ऐसी चुनी जाएँ कि उनका
अभिगम्यता संबंध स्वाभाविक रीतियों से प्रतिबंधित हो, तो प्रायः ऐसा अर्थविज्ञान
मिल जाता है जिसके लिए !!{derivation}-तंत्र सुदृढ़ और पूर्ण हो। उदाहरणतः
$\Log{S4}$ के अनुरूप सभी क्रिप्के !!{structure}s ऐसी हैं कि उनका अभिगम्यता
संबंध स्वतुल्य और संक्रामक है। $\Log{S5}$ के अनुरूप सभी क्रिप्के
!!{structure}s ऐसी हैं कि उनका अभिगम्यता संबंध {\em सार्वत्रिक} है, अर्थात् प्रत्येक जगत हर अन्य
जगत से अभिगम्य है; अतः $\Box !A$ तभी और केवल तभी सत्य है जब $!A$ प्रत्येक
जगत में सत्य हो।

\end{document}
